\chapter{Results}
\label{ch:results}

% Every number in this file was computed from the notebook (blocks 1-16) on:
%   - 0day _In the Wild_.xlsx, sheet "All"
%   - reserved_dates.csv (CVE List, pulled 8 October 2026)
%   - cvss_scores.csv (NVD via the fkie-cad mirror, pulled 8 October 2026)
% Two latency measures are shown side by side:
%   discovery-to-patch   = Date Patched - Date Discovered   (n = 132)
%   reservation-to-patch = Date Patched - CVE dateReserved  (n = 403)
% Which one is the headline is still to be agreed with the supervisor.
% Lines starting with "% ->" are the facts each paragraph has to state. Write the sentences yourself.

\section{Data Overview}
% -> 410 CVEs in the snapshot after cleaning
% -> discovery date known for 133 (32.4%); 132 of those also have a patch date (CVE-2024-7262 has none)
% -> reserved date found for all 410; 403 give a usable latency
%    (6 dropped because the ID was reserved after the patch date, 1 has no patch date)
% -> 18 CVEs have a reservation-to-patch latency of 0 (ID reserved on the patch day)
% -> medians: 9 days (discovery) vs 56 days (reservation); refer to the table
% -> both are right-skewed (mean 34.2 vs median 9; mean 80.7 vs median 56), so medians are reported throughout

\begin{table}[h]
    \centering
    \caption{Patch latency in days under both measures.}
    \label{tab:descriptives}
    \begin{tabular}{lrr}
        \toprule
        Statistic & Discovery-to-patch & Reservation-to-patch \\
        \midrule
        Observations             & 132 & 403 \\
        Median                   & 9 & 56 \\
        Interquartile range      & 3 -- 30.2 & 23.5 -- 114.5 \\
        90th percentile          & 85.5 & 186 \\
        Minimum                  & 1 & 0 \\
        Maximum                  & 834 & 426 \\
        Mean                     & 34.2 & 80.7 \\
        Patched within 7 days    & 48.5\% & 12.9\% \\
        Patched within 30 days   & 75.0\% & 28.5\% \\
        Patched within 90 days   & 90.2\% & 66.3\% \\
        \bottomrule
    \end{tabular}
\end{table}

\section{Proxy Validation}
% -> 133 CVEs have both a discovery and a reserved date
% -> mean absolute difference 55.7 days, standard deviation 88.3 days
% -> Spearman between the two dates: rho = 0.994 (high by construction: the dates span ten years)
% -> Spearman between the two latencies: rho = 0.438, p < 0.001, n = 132 (this is the one that matters)
% -> on those 132 CVEs the median latency is 9 days by discovery date and 57 by reserved date
% -> the ID was reserved before discovery in 83 cases (62%), on the same day in 1, after in 49
% -> median lead of reservation over discovery differs by vendor:
%    Microsoft 79 days, Adobe 31, Google 6, Apple 6
% -> and by period: 54 days up to 2021 (n = 90), -2 days from 2022 (n = 40);
%    from 2022, 75% are within a week of discovery
% -> consequence: the proxy adds a lead time that varies by vendor and by year, so it is not neutral for RQ1 and RQ2

\section{Normality}
% -> Shapiro-Wilk rejects normality for both measures (p < 0.001)
%    discovery: W = 0.371, skewness 7.38; reservation: W = 0.855, skewness 1.50
% -> consequence: all tests below are the non-parametric ones named in the methodology
% Two or three sentences.

\section{Patch Latency over Time (RQ1)}
% -> discovery measure (2015--2024): Spearman rho = -0.043, p = 0.628; Kendall tau = -0.026, p = 0.680. No trend.
%    yearly medians range from 3 (2022) to 37.5 (2024) with no direction
% -> reservation measure (2013--2026): Spearman rho = -0.440, p < 0.001; Kendall tau = -0.312, p < 0.001. Apparent decline.
% -> check on the 130 CVEs that have both measures, same year variable:
%    reservation-to-patch vs year: rho = -0.422, p < 0.001
%    discovery-to-patch vs year:   rho = 0.003, p = 0.969
%    reservation lead vs year:     rho = -0.552, p < 0.001
% -> so the decline comes from IDs being reserved closer to discovery in recent years, not from faster patching
% -> the hypothesis that latency increases with year is not supported under either measure
% -> caveats: few discovery-dated cases in 2017-2019; reservations from 2026 are cut short by the snapshot (median 3 (n = 14))

\begin{figure}[h]
    \centering
    % \includegraphics[width=0.48\textwidth]{figures/median_latency_per_year.png}
    % \includegraphics[width=0.48\textwidth]{figures/median_latency_per_year_reserved.png}
    \caption{Median patch latency per year, by discovery date (left) and by CVE reservation date (right).}
    \label{fig:latency-per-year}
\end{figure}

\section{Patch Latency across Vendors (RQ2)}
% -> Kruskal-Wallis across the four vendors:
%    discovery:   H = 36.68, p < 0.001 (n = 102)
%    reservation: H = 29.65, p < 0.001 (n = 337)
% -> robustness check, reservations from 2022 only, three vendors (Adobe has 2 cases):
%    H = 30.64, p < 0.001 (n = 157); medians Microsoft 42 (n = 65), Google 1 (n = 50), Apple 40 (n = 42)
% -> holds in all three views: vendors differ overall, and Google is faster than Microsoft
% -> does not hold across views: Apple's position; Adobe cannot be placed (mostly older cases)
% -> this is heterogeneity, not a ranking: vendor is confounded with product mix, release cadence and period

\begin{table}[h]
    \centering
    \caption{Patch latency in days per vendor under both measures.}
    \label{tab:vendors}
    \small
    \setlength{\tabcolsep}{4pt}
    \begin{tabular}{lrrrrrrrr}
        \toprule
        & \multicolumn{4}{c}{Discovery} & \multicolumn{4}{c}{Reservation} \\
        \cmidrule(lr){2-5} \cmidrule(lr){6-9}
        Vendor & $n$ & Median & IQR & Range & $n$ & Median & IQR & Range \\
        \midrule
        Google    & 53 & 4 & 2 -- 6 & 1 -- 132 & 79 & 45 & 1 -- 92 & 0 -- 426 \\
        Adobe     & 14 & 6 & 4.2 -- 8.5 & 1 -- 29 & 28 & 52.5 & 17.2 -- 95 & 5 -- 235 \\
        Apple     & 10 & 9.5 & 9 -- 10 & 6 -- 10 & 64 & 46.5 & 15.8 -- 91 & 4 -- 244 \\
        Microsoft & 25 & 24 & 15 -- 48 & 5 -- 182 & 166 & 62 & 40 -- 143.8 & 5 -- 379 \\
        \bottomrule
    \end{tabular}
\end{table}

\begin{table}[h]
    \centering
    \caption{Bonferroni-adjusted $p$-values of pairwise Mann-Whitney U tests between vendors. The first two columns are adjusted for six comparisons, the last for three.}
    \label{tab:pairwise}
    \begin{tabular}{lrrr}
        \toprule
        Pair & Discovery & Reservation & Reservation, 2022 on \\
        \midrule
        Microsoft -- Google & $<0.001$ & $<0.001$ & $<0.001$ \\
        Microsoft -- Apple  & 0.001 & $<0.001$ & 0.981 \\
        Microsoft -- Adobe  & $<0.001$ & 0.284 & -- \\
        Google -- Apple     & 0.027 & 0.245 & $<0.001$ \\
        Google -- Adobe     & 1.000 & 0.447 & -- \\
        Apple -- Adobe      & 0.116 & 1.000 & -- \\
        \bottomrule
    \end{tabular}
\end{table}

\begin{figure}[h]
    \centering
    % \includegraphics[width=0.48\textwidth]{figures/latency_by_vendor.png}
    % \includegraphics[width=0.48\textwidth]{figures/latency_by_vendor_reserved.png}
    \caption{Patch latency per vendor: discovery-to-patch on a logarithmic scale (left) and reservation-to-patch on a linear scale (right).}
    \label{fig:latency-by-vendor}
\end{figure}

\section{Patch Latency across Products (RQ3)}
% -> Google, Chrome vs Android: differs under both measures
%    discovery: Mann-Whitney U = 5, p < 0.001; reservation: Mann-Whitney U = 143, p < 0.001
%    also holds for reservations from 2022 only (Chrome 1 vs Android 96.5 days, p < 0.001)
% -> Microsoft, Windows vs Internet Explorer vs Office:
%    discovery: Kruskal-Wallis H = 1.06, p = 0.588 (groups of 3, cannot detect a difference)
%    reservation: Kruskal-Wallis H = 7.91, p = 0.019, driven by Windows vs Internet Explorer (adjusted p = 0.015)
%    but the products come from different periods (median reservation year: Internet Explorer 2016, Office 2016, Windows 2022);
%    restricted to reservations up to 2021: H = 0.99, p = 0.61 (medians: Internet Explorer 141, Office 120.5, Windows 124). No product difference.
% -> Apple, iOS vs WebKit: reservation: Mann-Whitney U = 361, p = 0.49. No difference. Discovery set too small (8 and 2).
% -> the hypothesis of a product difference within a vendor is supported for Google only

\begin{table}[h]
    \centering
    \caption{Patch latency in days per product under both measures.}
    \label{tab:products}
    \small
    \setlength{\tabcolsep}{4pt}
    \begin{tabular}{llrrrrrr}
        \toprule
        & & \multicolumn{2}{c}{Discovery} & \multicolumn{4}{c}{Reservation} \\
        \cmidrule(lr){3-4} \cmidrule(lr){5-8}
        Vendor & Product & $n$ & Median & $n$ & Median & IQR & Range \\
        \midrule
        Microsoft & Windows           & 13 & 27 & 98 & 60 & 40 -- 131.5 & 5 -- 379 \\
                  & Internet Explorer & 3 & 9 & 23 & 130 & 81 -- 212 & 11 -- 351 \\
                  & Office            & 3 & 14 & 18 & 66 & 42 -- 162.8 & 27 -- 242 \\
        \addlinespace
        Google    & Chrome            & 47 & 4 & 61 & 3 & 0 -- 71 & 0 -- 153 \\
                  & Android           & 6 & 116 & 18 & 115.5 & 76 -- 136.8 & 14 -- 426 \\
        \addlinespace
        Apple     & iOS               & 8 & 9.5 & 30 & 22 & 13.5 -- 101 & 4 -- 244 \\
                  & WebKit            & 2 & 9.5 & 27 & 49 & 20 -- 77 & 5 -- 240 \\
        \bottomrule
    \end{tabular}
\end{table}

\section{Patch Latency and Severity (RQ4)}
% -> CVSS v3 base score found for all 410 CVEs; 313 are NVD's own assessment, 97 come from the vendor or CISA
% -> Low has a single CVE and is left out of the tests
% -> discovery: Kruskal-Wallis H = 30.87, p < 0.001; Spearman rho = -0.648, p < 0.001
%    adjusted pairwise: Medium vs High p < 0.001; Medium vs Critical p < 0.001; High vs Critical p = 0.005
% -> reservation: Kruskal-Wallis H = 9.33, p = 0.009; Spearman rho = -0.216, p < 0.001
%    adjusted pairwise: Medium vs High p = 1.000; Medium vs Critical p = 0.033; High vs Critical p = 0.011
% -> direction matches the hypothesis under both measures: higher severity, shorter latency. Strong on discovery dates, weak on reserved dates.
% -> within vendors, discovery: Microsoft rho = -0.52 (p = 0.01, n = 25); Google rho = -0.52 (p < 0.001, n = 53); Apple rho = 0.19 (p = 0.59, n = 10); Adobe rho = 0.03 (p = 0.91, n = 14)
% -> within vendors, reservation: Microsoft rho = -0.02 (p = 0.78, n = 166); Google rho = -0.40 (p < 0.001, n = 79); Apple rho = -0.11 (p = 0.39, n = 64); Adobe rho = -0.08 (p = 0.67, n = 28)
% -> within Chrome alone, reservation: rho = -0.06, p = 0.64, n = 61 (the Google result reflects Chrome vs Android)
% -> caveats: 305 of 410 CVEs are High and 262 have a score of 7.8 or 8.8, so the scale barely varies;
%    bands are tied to vendors (discovery set: 8 of 21 Medium cases are Samsung, 20 of 25 Critical cases are Google or Adobe)

\begin{table}[h]
    \centering
    \caption{Patch latency in days per CVSS v3 severity band under both measures.}
    \label{tab:severity}
    \small
    \setlength{\tabcolsep}{4pt}
    \begin{tabular}{lrrrrrrr}
        \toprule
        & & \multicolumn{3}{c}{Discovery} & \multicolumn{3}{c}{Reservation} \\
        \cmidrule(lr){3-5} \cmidrule(lr){6-8}
        Band & CVEs & $n$ & Median & IQR & $n$ & Median & IQR \\
        \midrule
        Low      & 1 & 1 & 78 & -- & 1 & 39 & -- \\
        Medium   & 48 & 21 & 69 & 27 -- 121 & 47 & 60 & 40 -- 107 \\
        High     & 305 & 85 & 9 & 4 -- 24 & 300 & 57 & 28 -- 123.2 \\
        Critical & 56 & 25 & 4 & 3 -- 6 & 55 & 25 & 5.5 -- 85.5 \\
        \bottomrule
    \end{tabular}
\end{table}
